% !TeX root = ../navier-stokes-manuscript.tex
% ============================================================
% 2.1 Critical Height at a Finite Endpoint
% ============================================================

\subsection{Critical Height at a Finite Endpoint}\label{sub:CriticalHeightAtAFiniteEndpoint}

Fix a rapidly decaying datum \(u_0\) on \(\R^3\), a viscosity \(\nu>0\), and the maximal classical solution \((u,p)\) they generate on \([0,T_*)\).
Assume \(T_*<\infty\).
The dilation and the decaying pressure normalization used below are whole-space statements.
On \(\mathbb T^3\), discrete frequencies replace this dilation and require a separate high-frequency or localization argument.

The solution is smooth at every earlier time, so any singular behaviour must be approached through a sequence of regular states as \(t\uparrow T_*\).
Kinetic energy controls the total \(L^2\) size of those states.
That \(L^2\) ceiling remains compatible with rescaled states at progressively shorter length scales.
Navier--Stokes scaling singles out the Sobolev height that stays unchanged under exactly that compression.

If \((u,p)\) solves the Navier--Stokes equations with viscosity \(\nu>0\), then for every \(\lambda>0\),
\[
  u_\lambda(x,t)
  :=
  \lambda u(\lambda x,\lambda^2t),
  \qquad
  p_\lambda(x,t)
  :=
  \lambda^2p(\lambda x,\lambda^2t)
\]
solves the same equation with the same viscosity on the rescaled time interval.
The spatial scale is divided by \(\lambda\), while the corresponding time scale is divided by \(\lambda^2\).
This is the parabolic relation already present in \(\partial_t-\nu\Delta\).

Let \(\Lambda=(-\Delta)^{1/2}\), and write
\[
  \norm{v}_{\dot H^s}^2
  :=
  \int_{\R^3}
  |\xi|^{2s}|\widehat v(\xi)|^2
  \,d\xi.
\]
The Fourier transform of the rescaled velocity is
\[
  \widehat{u_\lambda}(\xi,t)
  =
  \lambda^{-2}
  \widehat u(\xi/\lambda,\lambda^2t).
\]
Substituting this expression and setting \(\eta=\xi/\lambda\) gives
\begin{align*}
  \norm{u_\lambda(t)}_{\dot H^s}^2
  &={}
  \int_{\R^3}
  |\xi|^{2s}\lambda^{-4}
  \left|\widehat u(\xi/\lambda,\lambda^2t)\right|^2
  \,d\xi
  \\
  &={}
  \lambda^{2s-1}
  \int_{\R^3}
  |\eta|^{2s}
  \left|\widehat u(\eta,\lambda^2t)\right|^2
  \,d\eta.
\end{align*}
Hence
\[
  \norm{u_\lambda(t)}_{\dot H^s}
  =
  \lambda^{s-1/2}
  \norm{u(\lambda^2t)}_{\dot H^s}.
\]
The exponent vanishes at \(s=1/2\).
This gives the critical height
\[
  H(t)
  :=
  \frac12\norm{u(t)}_{\dot H^{1/2}}^2
  =
  \frac12\norm{\Lambda^{1/2}u(t)}_2^2.
\]
The scaling can compress a smooth configuration to a finer spatial length and a shorter time while preserving the value of \(H\).

At the kinetic-energy level, the same calculation gives
\[
  \norm{u_\lambda(t)}_2^2
  =
  \lambda^{-1}
  \norm{u(\lambda^2t)}_2^2.
\]
Successive rescalings can therefore lower the kinetic-energy scale while leaving the critical height unchanged.
This is the precise opening through which a finite-energy estimate can fail to see a concentration at ever finer scales.

The finite endpoint gives the scaling a second consequence.
Choose a geometric sequence of lengths
\[
  \ell_k
  :=
  2^{-k}\ell_0.
\]
The viscous time associated with \(\ell_k\) has parabolic size
\[
  \tau_k
  :=
  \frac{c\ell_k^2}{\nu}
  =
  \frac{c\ell_0^2}{\nu}4^{-k},
\]
where \(c>0\) fixes the chosen normalization.
These times have a finite sum,
\[
  \sum_{k=0}^{\infty}\tau_k
  =
  \frac{4c\ell_0^2}{3\nu}
  <
  \infty.
\]
Parabolic timing therefore permits an abstract Zeno schedule in which infinitely many progressively finer episodes fit before a finite time.
The calculation establishes a dimensionally consistent schedule.
Realizing that schedule would require the same Navier--Stokes history to generate every episode while transport, pressure, incompressibility, and viscosity continue to act together.

The critical height tells us which spatial measurement can remain visible through such a schedule.
The speed at which the field moves from one regular state to the next is still unknown.
Finite time makes that response question exact: an unbounded rung can be reached only through the accumulated action of the rung above it.
